% Provenance Is Not Validity: The Vacuity of Cryptographically Verified Observation
% Critique paper skeleton. House style: plain declarative, no em-dashes in prose.
% Build: pdflatex poso ; bibtex poso ; pdflatex poso ; pdflatex poso
% Class is a starting point. Swap to a systems / metascience / astro-methods template as needed.
\documentclass[conference]{IEEEtran}

\usepackage[T1]{fontenc}
\usepackage{lmodern}
\usepackage{amsmath,amssymb}
\usepackage{booktabs}
\usepackage{graphicx}
\usepackage{url}
\usepackage[hidelinks]{hyperref}

\begin{document}

\title{Provenance Is Not Validity:\\ The Vacuity of Cryptographically Verified Observation}

\author{\IEEEauthorblockN{Andrew H. Bond}
\IEEEauthorblockA{San Jos\'e State University\\
San Jos\'e, CA, USA\\
% TODO: contact email (institutional preferred)
Email: [TODO]}}

\maketitle
\thispagestyle{empty}
\pagestyle{empty}

\begin{abstract}
A growing class of systems anchors scientific measurements on a blockchain to make them
trustworthy. SkyMapper L1 and its Proof of Space Observation are a current example: telescope
observations are signed at the source, time-stamped, and recorded on a distributed ledger, and
are said to become provably reliable and reproducible. We show that this claim conflates two
guarantees that must be kept apart. Cryptographic anchoring establishes provenance, the
authenticity, integrity, and time of a record. It does not, and by construction cannot, establish
validity, whether the instrument was calibrated, whether it pointed where it claims, and whether
the observation supports a downstream scientific read. A miscalibrated telescope's output signs
and anchors exactly as cleanly as a correct one, so the ledger clears every record with equal
confidence regardless of scientific worth. We formalize this through consumer-relative observation
and characterize such a system as a false-clear mechanism: it certifies the source, not the reader
who must use the data. We then supply the missing layer. A validity attestation re-evaluated at
use time, a measured false-clear rate for the network, and re-observation as the consumer test
convert provenance into evidence. On real Breakthrough Listen L-band data, an ON-OFF cadence
consumer rejects roughly 99.95\% of raw single-scan detections as interference, a discrimination
no signature provides. For trusted observational datasets provenance is necessary and not
sufficient, and the validity layer, not the ledger, is where trust is earned.
\end{abstract}

\begin{IEEEkeywords}
scientific data integrity; provenance versus validity; consumer-relative observation; false-clear
rate; technosignature verification; blockchain critique; trusted datasets
\end{IEEEkeywords}

% ===================================================================================
\section{Introduction}
\label{sec:intro}
% Claim in plain words, then the two or three findings as prose, then credit what is real.
Observational science is acquiring a new trust infrastructure. Systems now sign a measurement at
the sensor, time-stamp it, and anchor it on a public ledger, and present the result as data that
can be trusted without trusting the operator. SkyMapper L1 and its Proof of Space Observation
(POSO) are a current and well-funded instance, contributing live telescope data across many sites
and naming reliable, tamper-proof, and reproducible datasets as the outcome~\cite{skymapper}.

We make one argument. Cryptographic anchoring proves that a record is authentic, unaltered, and
old. It does not prove that the record is right. Those are different guarantees, and only the
first is a property of a signature. The second is a property of the instrument, the calibration,
and the downstream use, and no amount of anchoring reaches it. A ledger that anchors a
miscalibrated telescope's output has proven the wrong thing well.

We are careful about what the cryptography does earn. For attribution, where the requirement is
non-repudiable chain of custody, provenance is exactly right. For pre-registration, where a
timestamp that provably predates a claim resists after-the-fact story fitting, provenance is again
the right tool. The paper does not dispute these. It disputes the slide from provenance to
validity that the reliability and reproducibility framing performs.

The contribution is threefold. First, we separate provenance from validity precisely and locate
the conflation in the class of onchain-observation systems, with POSO as the motivating case
(Sections~\ref{sec:background} and~\ref{sec:distinction}). Second, we characterize such a system,
through consumer-relative observation, as a false-clear mechanism, and define a measurable
false-clear rate for it (Section~\ref{sec:falseclear}). Third, we give the validity layer that the
provenance layer lacks, and we demonstrate the load-bearing part of it, a re-observation consumer,
on real Breakthrough Listen data (Sections~\ref{sec:layer} and~\ref{sec:evidence}).

% ===================================================================================
\section{Background: onchain observation}
\label{sec:background}
% What POSO and its class actually do. Keep this factual and sourced.
% TODO: replace paraphrase with sourced specifics once the primary description is accessible.
The pattern is consistent across the class. A sensor holds a signing key. Each observation, or a
digest of it, is signed at the point of capture, associated with a time, and written to or
anchored on a distributed ledger, with the bulk data held in decentralized or content-addressed
storage~\cite{skymapper}. The stated goals are tamper-evidence, non-repudiation, and an auditable
history that any party can check without trusting a central operator.

We treat POSO as a pure-provenance mechanism, which is the design its public claims describe. This
is the strongest reading for the system, because it takes the mechanism at face value. If a later
specification adds instrument-level checks, the critique in Section~\ref{sec:distinction} narrows
from the design to the claim, and the constructive layer of Section~\ref{sec:layer} still applies.
% TODO: if POSO's spec becomes available, state here exactly what, if anything, it does beyond
% sign/timestamp/anchor, and quote the reliability/reproducibility claims verbatim.

Cryptographic timestamping and tamper-evident logs are mature and well understood outside the
ledger setting~\cite{rfc3161, merkle}. A blockchain is one way to obtain tamper-evident ordering
and open verifiability. Whether it is the most economical way is not our question. Our question is
what the ordering and the signatures are \emph{about}.

% ===================================================================================
\section{Two guarantees, kept apart}
\label{sec:distinction}
Anchoring establishes three facts, and all three concern the bytes. Authenticity: a known key
signed this record. Integrity: the record is unchanged since it was signed. Existence at a time:
the record existed by a given timestamp. Call the conjunction \emph{provenance}.

Validity is a different property. It asks whether the instrument was calibrated within its
known-good envelope, whether it pointed where the metadata claims, whether the exposure was free of
saturation, artifacts, and interference, and whether the recorded observation supports the specific
read a downstream consumer intends. Call this \emph{validity}.

The two are independent. A saturated sensor, a mispointed mount, a hot pixel, or an uncalibrated
gain produces a record that signs and anchors exactly as cleanly as a correct observation. The
signature is a function of the bytes and the key. It is invariant to whether the bytes describe a
good measurement. Therefore anchoring cannot separate a valid observation from an invalid one, and
a system built only on anchoring clears every record with equal confidence regardless of its
scientific worth.

The reliability and reproducibility claims cross this line. Reliability is a validity property of
the instrument. Reproducibility is a property of independent re-observation and re-derivation, and
a signature never re-observes anything. Only tamper-evidence, the middle term, is a provenance
property that the cryptography actually delivers.

% ===================================================================================
\section{A false-clear characterization}
\label{sec:falseclear}
% Formalize via consumer-relative observation. Keep the math light: 2-4 numbered equations max.
Observation Theory holds that an observation is meaningful only relative to the consumer that
reads it, and that a certificate is worth only what it lets that consumer verify~\cite{ot}. Let a
record $x$ carry a certificate $C(x)$, and let a downstream consumer apply a validity test
$V(x) \in \{0,1\}$, where $V(x)=1$ means the record supports the consumer's read. A certificate is
\emph{consumer-sound} when it predicts $V$:
\begin{equation}
C(x)=1 \;\Rightarrow\; V(x)=1 .
\end{equation}
A provenance certificate $C_{\mathrm{prov}}$ attests only that the source signed $x$ and that $x$
is unaltered. It is by construction independent of $V$, so
\begin{equation}
\Pr\!\left[V(x)=0 \,\middle|\, C_{\mathrm{prov}}(x)=1\right] \;=\; \Pr\!\left[V(x)=0\right],
\end{equation}
the base rate of invalid records. We call this conditional probability the \emph{false-clear rate}
of the certificate. A provenance certificate does not reduce it below the base rate. When invalid
records are common, which for raw observational data they are, the false-clear rate is high and
the certificate is, in the theory's sense, vacuous. The failure is not weak cryptography. The
cryptography is sound. The failure is that $C_{\mathrm{prov}}$ certifies the source, not the reader
whose test is $V$.

% ===================================================================================
\section{The validity layer}
\label{sec:layer}
Four moves make a certificate consumer-sound, and each maps onto machinery that already exists.

\textbf{Validity attestation, re-evaluated at use time.} The certificate should attest what a
consumer can validly extract, bound to calibration state and instrument envelope, and it should be
re-evaluated when the data are used rather than fixed at capture. This is the validity-attestation
class and use-time re-evaluation of IEEE P3787~\cite{p3787}.
% TODO: cite the exact P3787 clause for the validity-attestation class.

\textbf{A measured false-clear rate for the network.} Of the records a system clears, the fraction
that fail a real consumer test $V$ is the honest worth of its certificate. Reporting this rate
turns a trust claim into a measurement, and a high rate is a falsification.

\textbf{Re-observation as the consumer test.} For observational data the most consumer-relevant
$V$ is persistence: a signal that recurs under independent looks, and is absent where it should be
absent, is supported in a way a single signed frame can never be. This is the ON-OFF cadence logic
of technosignature search, promoted from a pipeline convenience to the definition of the
certificate.

\textbf{An in-band witness.} The witness should live in the observation's own read path, so that it
is coupled to the instrument's fidelity, rather than a separate anchor that certifies the record's
history while decoupled from its measurement quality.

% ===================================================================================
\section{Evidence: a measured consumer test}
\label{sec:evidence}
% The empirical anchor. Numbers here come from the graded run; do not quote shakedown values.
We demonstrate the re-observation consumer on real data, to show that a validity test does work
that a signature cannot. We ran the standard narrowband technosignature search on a real six-scan
ON-OFF cadence of a nearby star, observed at L-band with the Green Bank Telescope and released
through Breakthrough Listen~\cite{bl-opendata, turboseti}. The raw search flags every narrowband
hit in each scan. The cadence consumer keeps only a hit that appears in the on-source scans and is
absent from the off-source pointings.

Across the three on-source scans the raw search returned $6651$ hits, almost all of them
attributable to Global Positioning System and satellite interference near $1575$~MHz. The cadence
consumer promoted $3$ of them to candidates. The false-clear rate of the raw single-scan detector,
relative to the cadence consumer, is therefore $1 - 3/6651 \approx 0.9995$. A signature on any of
these frames would have certified all $6651$ as authentic and unaltered, and would have separated
none of them. The discrimination comes from re-observation, not from provenance.
% TODO: Fig.~\ref{fig:cadence}: waterfall of the ON and OFF scans with the surviving candidate,
% built from cadence_result.json (see /archive/seti on the compute host).
% TODO: state the target designation, MJD, and coarse-channel range in a parameters table.

We report this as a mechanism demonstration, not as a benchmark of POSO, which we did not run. It
establishes that a consumer test yields a large, measurable separation on real observational data,
which is precisely the quantity a provenance certificate leaves at its base rate.

% ===================================================================================
\section{Related work}
\label{sec:related}
% Credit prior art precisely. End with the "we offer the measurement" move.
% TODO: verify every entry against a real record (authors, venue, year) before submission.
Blockchain proposals for scientific data provenance and reproducibility are numerous, and critical
appraisals of them exist~\cite{blockchain-science, blockchain-critique}. Tamper-evident logging and
trusted timestamping predate the ledger framing~\cite{rfc3161, merkle}. The distinction between
data provenance and data quality is long established in the databases and metrology
literatures~\cite{provenance-survey, dataquality}, and our contribution is to carry it into the
onchain-observation setting and to attach a measurable rate to it. Consumer-relative observation
and the validity-attestation program provide the frame~\cite{ot, p3787}. We do not propose a better
ledger. We propose the measurement that shows what a ledger does not do.

% ===================================================================================
\section{Discussion and scope}
\label{sec:scope}
% Limits stated flatly. Negatives disclosed.
The critique is narrow and precise. Provenance is necessary for trusted observational datasets. It
is not sufficient. We do not claim that anchoring is worthless. For attribution and for verifiable
pre-registration it is the right tool, and we say so in Section~\ref{sec:intro}.

We read SkyMapper's public description and treat POSO as a pure-provenance mechanism. If its
specification incorporates instrument-level validity checks, the critique narrows from the design
to the claim, and the constructive layer still applies. Our empirical demonstration uses one
cadence, one frequency sub-band, and one narrowband consumer. It is a demonstration that a consumer
test separates real data, not a survey-scale statistic, and we do not generalize it beyond that.

% ===================================================================================
\section{Conclusion}
\label{sec:conclusion}
Onchain anchoring proves that an observation is authentic. It does not prove that the observation
is right, and for science whether it is right is the whole question. The validity layer, not the
ledger, is where trust is earned. We have separated the two guarantees, characterized a
provenance-only system as a false-clear mechanism with a base-rate false-clear rate, given the
layer that closes the gap, and shown on real data that a consumer test delivers the separation a
signature cannot.

% ===================================================================================
\bibliographystyle{IEEEtran}
\bibliography{refs}

\end{document}
